\documentclass[UTF8,a4paper,11pt]{ctexart}
\usepackage[a4paper,margin=2.3cm]{geometry}
\usepackage{amsmath,amssymb,mathtools}
\usepackage{enumitem}
\usepackage[hidelinks]{hyperref}
\usepackage{fancyhdr}
\setlist[itemize]{leftmargin=2.2em,itemsep=0.25em,topsep=0.35em}
\setlength{\parindent}{2em}\setlength{\parskip}{0.35em}
\pagestyle{fancy}\fancyhf{}\fancyhead[L]{离散数学与结构（I）}\fancyhead[R]{Lec 17}\fancyfoot[C]{\thepage}\setlength{\headheight}{14pt}
\title{熵、TV 距离与 KL 散度}\author{zhangjunshuo}\date{2026 年 11 月}
\begin{document}
\maketitle
\begin{quote}
本讲根据 2026 年 11 月 12 日的两页手写稿和一段课堂转写整理。上一份笔记是 Lec16（11 月 10 日），本次按顺序记为 \textbf{Lec17}。转写开头仍在讲上一讲的 PGF 例子（等待连续两个正面），之后进入熵与信息论；整理版把它作为第 1 节的收尾。语音识别对术语和公式错误很多，公式以手写稿为准，转写用于补课堂解释；不一致处见总览中的勘误。

本讲主线：PGF 收尾 $\to$ 熵与二元熵 $\to$ 条件熵/联合熵/链式法则 $\to$ TV 距离与耦合 $\to$ KL 散度与条件 KL。

每个分节开头都先交代本节用到的记号；所有符号在首次使用处给出定义。
\end{quote}
\tableofcontents\clearpage

\section{PGF 收尾、熵的定义与基本性质}


这一讲先用一个“等待连续两个正面”的例子把概率生成函数（PGF）收尾，然后正式进入信息论的主角：\textbf{熵}。熵衡量的是一个随机变量平均能提供多少信息，或者“描述它平均需要多少比特”。

\textbf{本节用到的记号：} $p=P(H)$，$q=1-p$；$X$ 为等待时间；$P_X$ 为分布；$\operatorname{supp}(P_X)=\{x:P_X(x)>0\}$ 为支撑集；$\log$ 的底数暂不固定，$H$ 随底数只差一个常数因子。

\subsection{上一讲收尾：等待连续两个正面的 PGF}


设每次抛硬币独立，正面概率为 $p$，反面概率为 $q=1-p$。令 $X$ 为\textbf{第一次出现连续两个正面（HH）所需的抛掷次数}。求 $X$ 的 PGF 和期望。

\textbf{状态法.} 记

\begin{itemize}
\item $A(z)$：当前没有“上一个正面”时，从此刻到首次 HH 的 PGF；
\item $B(z)$：当前刚抛出一个正面时，从此刻到首次 HH 的 PGF。
\end{itemize}


从状态 $A$：

\begin{itemize}
\item 抛反面（概率 $q$）：消耗 $1$ 次，回到 $A$，贡献 $qzA(z)$；
\item 抛正面（概率 $p$）：消耗 $1$ 次，进入 $B$，贡献 $pzB(z)$。
\end{itemize}


所以

\[
A(z)=z\bigl(qA(z)+pB(z)\bigr).
\]

从状态 $B$：

\begin{itemize}
\item 再抛正面（概率 $p$）：立即完成，贡献 $pz$；
\item 抛反面（概率 $q$）：消耗 $1$ 次，回到 $A$，贡献 $qzA(z)$。
\end{itemize}


所以

\[
B(z)=z\bigl(p+qA(z)\bigr).
\]

解这个方程组。由第一式

\[
A(z)(1-qz)=pzB(z),
\]

代入第二式：

\[
A(z)\bigl(1-qz-pqz^2\bigr)=p^2z^2.
\]

因此

\[
\boxed{
G_X(z)=A(z)
=\frac{p^2z^2}{1-qz-pqz^2}
=\frac{p^2z^2}{1-(1-p)z-p(1-p)z^2}
}.
\]

\textbf{期望.} 对 $z$ 求导并在 $z=1$ 取值，可得

\[
E[X]=G_X'(1)=\frac{1+p}{p^2}.
\]

检验：公平硬币 $p=1/2$ 时 $E[X]=6$，与“平均 6 次抛出 HH”的经典结论一致。

\textbf{直观.} 这个例子说明 PGF 不仅能读分布，还能把“带吸收状态的等待时间”通过一次方程解出来；后面信息论里遇到“平均需要多少符号/比特”时，也会反复使用“先写递推/方程，再读系数或导数”的思路。

\subsection{熵的定义}


设 $X$ 是离散随机变量，分布为 $P_X$。定义 $X$ 的熵为

\[
H(X)=H(P_X)
=\sum_{x}P_X(x)\log\frac{1}{P_X(x)}
=E\left[\log\frac{1}{P_X(X)}\right].
\]

约定 $0\log\frac10=0$；求和只在支撑集上进行，或把零概率项按此约定忽略。

\textbf{log 的底数.} 底数不固定，因为换底只差一个常数：

\begin{itemize}
\item 底 $2$：单位是 bit；
\item 底 $e$：单位是 nat；
\item 底 $256$：单位是 byte。
\end{itemize}


写 $\exp(\cdot)$ 时，指数函数的底要与 $\log$ 的底一致，否则公式会差常数。

\textbf{直观.} $-\log P_X(x)$ 是“看到 $x$ 时的惊讶程度/信息量”：概率越小，信息量越大；$H(X)$ 是它的平均值。

\subsection{二元熵}


最重要的例子是 Bernoulli 分布 $X\sim\mathrm{Bern}(p)$。记

\[
h(p)=H(\mathrm{Bern}(p))
=p\log\frac1p+(1-p)\log\frac1{1-p}.
\]

性质：

\begin{itemize}
\item $h(0)=h(1)=0$；
\item $h$ 在 $p=1/2$ 处最大，最大值为 $\log2$；
\item $h$ 关于 $p=1/2$ 对称：$h(p)=h(1-p)$；
\item $h$ 在 $[0,1]$ 上是\textbf{凹函数}。
\end{itemize}


若用 $\log_2$，则 $h(1/2)=1$ bit；这解释了“一个公平硬币正好提供 1 bit 信息”。

\subsection{熵的基本性质}


\textbf{非负性.} $H(X)\ge0$，且

\[
H(X)=0
\iff X\text{ 是常数（几乎必然）}.
\]

因为每一项 $P_X(x)\log\frac1{P_X(x)}\ge0$，且只有 $P_X(x)=1$ 时为零。

\textbf{上界.} 若 $|\operatorname{supp}(P_X)|=m$，则

\[
H(X)\le\log m,
\]

等号当且仅当 $X$ 在支撑集上均匀分布。

\textbf{证明（Jensen）.} $\log$ 是凹函数，所以

\[
\begin{aligned}
H(X)
&=\sum_xP_X(x)\log\frac1{P_X(x)}\\
&\le\log\left(\sum_xP_X(x)\cdot\frac1{P_X(x)}\right)\\
&=\log m.
\end{aligned}
\]

这里用了凹函数的 Jensen 形式：函数值的加权平均不超过加权平均的函数值，即 $\sum p_i\log y_i\le\log\sum p_i y_i$（当 $\sum p_i=1$）；取 $y_i=1/p_i$ 即得。

\textbf{函数处理不会增加熵.} 对任意函数 $f$，

\[
H(f(X))\le H(X),
\]

等号当 $f$ 在支撑集上是单射（可逆）时成立。直观：对 $X$ 做确定性处理，只会丢失信息，不会创造信息。

\textbf{凹性.} 若 $P=\lambda P_0+(1-\lambda)P_1$ 是两个分布的混合（$\lambda\in[0,1]$），则

\[
H(P)\ge\lambda H(P_0)+(1-\lambda)H(P_1).
\]

也就是说，熵作为分布的函数是凹的；混合分布比“各自熵的混合”更不确定。证明见下一节的条件熵（用 $H(X)\ge H(X\mid Z)$）。

\subsection{熵为什么自然：重复编码的直观}


若 $X$ 的熵是 $H$ 比特，独立重复 $n$ 次得到 $X_1,\dots,X_n$，则大概只需要 $nH$ 个比特就能以很小错误率描述这串结果，而不需要 $n\log|\mathcal X|$ 个比特。

例如 $X\sim\mathrm{Bern}(0.1)$，单次熵约 $0.469$ bit；重复 $10000$ 次，大约 $4690$ bit 就够，而不是 $10000$ bit。剩下的“小概率例外”可以用很小的错误概率处理。这正是后续压缩/信道编码讨论的起点。

\section{条件熵、联合熵与链式法则}


熵描述单个随机变量的不确定性；条件熵描述“已知 $Y$ 之后，$X$ 还剩多少不确定性”。本节的核心公式是链式法则

\[
H(X,Y)=H(Y)+H(X\mid Y),
\]

以及“条件不会增加熵”：

\[
H(X\mid Y)\le H(X).
\]

\textbf{本节用到的记号：} $P_{X,Y}(x,y)=P(X=x,Y=y)$；$P_{X\mid Y}(x\mid y)=P(X=x\mid Y=y)$；$H(X\mid Y=y)$ 为给定 $Y=y$ 的条件熵；$H(X\mid Y)$ 为对这些条件熵按 $Y$ 的分布取平均。

\subsection{条件熵的定义}


给定 $Y=y$，条件分布为 $P_{X\mid Y=y}$，其熵为

\[
H(X\mid Y=y)
=\sum_xP_{X\mid Y}(x\mid y)\log\frac1{P_{X\mid Y}(x\mid y)}.
\]

再对 $Y$ 的取值按 $P_Y(y)$ 加权平均，得到\textbf{条件熵}：

\[
H(X\mid Y)
=\sum_yP_Y(y)H(X\mid Y=y).
\]

\textbf{注意.} $H(X\mid Y=y)$ 是一个数；$H(X\mid Y)$ 是一个数（不是随机变量）——它已经把 $Y$ 平均掉了。若把 $H(X\mid Y)$ 写成随机变量，应先定义 $H(X\mid Y)(\omega)$，再取期望；本课程不需要这样写。

\textbf{特例.} 若 $Y$ 是常数，则 $H(X\mid Y)=H(X)$；若 $Y=X$，则 $H(X\mid Y)=0$。

\subsection{联合熵}


两个随机变量的联合熵为

\[
H(X,Y)
=\sum_{x,y}P_{X,Y}(x,y)\log\frac1{P_{X,Y}(x,y)}.
\]

它衡量的是“同时观察 $X,Y$ 时总的平均不确定性”。

\subsection{链式法则}


\textbf{定理.}

\[
\boxed{
H(X,Y)=H(Y)+H(X\mid Y)
}
\]

对称地也有 $H(X,Y)=H(X)+H(Y\mid X)$。

\textbf{证明.} 把联合分布写成 $P_{X,Y}(x,y)=P_Y(y)P_{X\mid Y}(x\mid y)$：

\[
\begin{aligned}
H(X,Y)
&=\sum_{x,y}P_{X,Y}(x,y)\log\frac1{P_{X,Y}(x,y)}\\
&=\sum_{x,y}P_{X,Y}(x,y)\left[
\log\frac1{P_Y(y)}+\log\frac1{P_{X\mid Y}(x\mid y)}
\right]\\
&=\sum_yP_Y(y)\log\frac1{P_Y(y)}
+\sum_{x,y}P_{X,Y}(x,y)\log\frac1{P_{X\mid Y}(x\mid y)}\\
&=H(Y)+H(X\mid Y).
\end{aligned}
\]

最后一项正是条件熵的定义。

\textbf{多变量链式法则.} 对 $X_1,\dots,X_n$，

\[
H(X_1,\dots,X_n)
=H(X_1)+H(X_2\mid X_1)+\cdots
+H(X_n\mid X_1,\dots,X_{n-1}).
\]

证明：反复用两变量链式法则，把最后一个变量单独提出来：

\[
H(X_1,\dots,X_n)=H(X_1,\dots,X_{n-1})+H(X_n\mid X_1,\dots,X_{n-1}),
\]

再对前 $n-1$ 个变量归纳。

\subsection{条件熵的非负性与“条件不增加熵”}


\textbf{非负性.}

\[
H(X\mid Y)\ge0.
\]

因为它是若干条件分布的熵的加权平均，而每个熵都非负。

\textbf{条件不增加熵.}

\[
\boxed{
H(X\mid Y)\le H(X)
}
\]

等号成立当且仅当 $X,Y$ 独立。

\textbf{证明（用 KL/对数和）.} 计算差值：

\[
\begin{aligned}
H(X)-H(X\mid Y)
&=\sum_xP_X(x)\log\frac1{P_X(x)}
-\sum_{x,y}P_{X,Y}(x,y)\log\frac1{P_{X\mid Y}(x\mid y)}\\
&=\sum_{x,y}P_{X,Y}(x,y)\log\frac{P_{X\mid Y}(x\mid y)}{P_X(x)}\\
&=\sum_{x,y}P_{X,Y}(x,y)\log\frac{P_{X,Y}(x,y)}{P_X(x)P_Y(y)}.
\end{aligned}
\]

这是 $P_{X,Y}$ 与“乘积分布”$P_XP_Y$ 之间的 KL 散度（见第 4 节）。由 KL 非负性可得差值 $\ge0$；等号当且仅当 $P_{X,Y}=P_XP_Y$，即 $X,Y$ 独立。

（如果还没学 KL，也可以直接用 Jensen：$\log$ 凹，$\sum p\log y\le\log\sum py$。）

\textbf{联合熵的上界.} 由链式法则和条件熵不等式，

\[
H(X,Y)=H(X)+H(Y\mid X)\le H(X)+H(Y),
\]

等号当且仅当 $X,Y$ 独立。多变量情形同样有

\[
H(X_1,\dots,X_n)\le\sum_iH(X_i).
\]

\subsection{条件熵为零的含义}


\[
H(X\mid Y)=0
\]

当且仅当：给定 $Y$ 后，$X$ 没有不确定性，即 $X$ 是 $Y$ 的函数（存在可测函数 $f$ 使 $X=f(Y)$ 几乎必然）。

这与“$H(X,Y)=H(Y)$”等价：联合的不确定性全在 $Y$ 上。

\subsection{熵的凹性（补证）}


上一节说 $H$ 是分布的凹函数。这里用条件熵给一个干净的证明。

设 $P=\lambda P_0+(1-\lambda)P_1$，$\lambda\in[0,1]$。构造随机变量 $Z\sim\mathrm{Bern}(\lambda)$，并令

\[
X\mid Z=0\sim P_0,\qquad X\mid Z=1\sim P_1.
\]

则 $X$ 的边缘分布正好是 $P$。由链式法则，

\[
H(X,Z)=H(X)+H(Z\mid X)=H(Z)+H(X\mid Z).
\]

其中：

\begin{itemize}
\item $H(Z)=h(\lambda)$；
\item $H(X\mid Z)=\lambda H(P_0)+(1-\lambda)H(P_1)$；
\item 由“条件不增加熵”，$H(Z\mid X)\le H(Z)$。
\end{itemize}


于是

\[
\begin{aligned}
H(P)=H(X)
&=H(Z)+H(X\mid Z)-H(Z\mid X)\\
&\ge H(X\mid Z)\\
&=\lambda H(P_0)+(1-\lambda)H(P_1).
\end{aligned}
\]

这正是凹性。等号当且仅当 $H(Z\mid X)=H(Z)$；在混合模型中，这通常要求 $P_0=P_1$（否则知道 $X$ 会改变 $Z$ 的后验）。

\textbf{结论.}

\[
H(\lambda P_0+(1-\lambda)P_1)
\ge\lambda H(P_0)+(1-\lambda)H(P_1).
\]

\subsection{常见错误}


\begin{itemize}
\item 把 $H(X\mid Y)$ 与 $H(X\mid Y=y)$ 混为一谈：前者是平均后的数。
\item 把链式法则写成 $H(X,Y)=H(X)+H(Y)$；只有独立时才能这样拆。
\item 认为 $H(X\mid Y)\le H(X)$ 的等号总是成立；等号要求独立。
\item 忘记“条件熵”是“平均条件熵”，不是“某个 $y$ 下的条件熵”。
\end{itemize}

\section{统计距离（TV 距离）}


比较两个概率分布有多“像”，最基础的工具是\textbf{统计距离}，也叫\textbf{全变差距离}（total variation distance, TV distance）。它的定义非常直观：把两个分布逐点相减，取绝对值再求和，最后除以 $2$。

\textbf{本节用到的记号：} $P,Q$ 是同一可数集合 $\mathcal X$ 上的两个概率分布；$P(x),Q(x)$ 为概率质量函数；$\Delta_{\mathrm{TV}}(P,Q)$ 为 TV 距离；$\operatorname{supp}(P)$ 为支撑集。

\subsection{定义与等价形式}


\[
\boxed{
\Delta_{\mathrm{TV}}(P,Q)
=\frac12\sum_{x\in\mathcal X}|P(x)-Q(x)|
}.
\]

它恰好是 $P-Q$ 的 $L^1$ 范数的一半。因为两个分布的总质量都为 $1$，所以：

\[
\Delta_{\mathrm{TV}}(P,Q)
=\sum_{x:P(x)>Q(x)}(P(x)-Q(x))
=\sum_{x:Q(x)>P(x)}(Q(x)-P(x)).
\]

也可以写成

\[
\Delta_{\mathrm{TV}}(P,Q)
=\sum_x\max\{0,P(x)-Q(x)\}.
\]

以及

\[
\Delta_{\mathrm{TV}}(P,Q)
=1-\sum_x\min\{P(x),Q(x)\}.
\]

最后一条的直观：两个分布“重叠”的质量是 $\sum_x\min(P(x),Q(x))$，剩下的部分就是它们不同的质量；由于总质量各为 $1$，不同部分两边各占一半，所以除以 $2$ 正好是 TV。

\textbf{例（Bernoulli）.} 若 $P=\mathrm{Bern}(p)$、$Q=\mathrm{Bern}(q)$，则

\[
\Delta_{\mathrm{TV}}(P,Q)=|p-q|.
\]

\textbf{基本性质.}

\begin{itemize}
\item $0\le\Delta_{\mathrm{TV}}(P,Q)\le1$；
\item 对称：$\Delta_{\mathrm{TV}}(P,Q)=\Delta_{\mathrm{TV}}(Q,P)$；
\item 同一分布距离为 $0$：$\Delta_{\mathrm{TV}}(P,P)=0$；
\item 满足三角不等式：$\Delta_{\mathrm{TV}}(P,R)\le\Delta_{\mathrm{TV}}(P,Q)+\Delta_{\mathrm{TV}}(Q,R)$。
\end{itemize}


所以 TV 是一个真正的距离（度量）。

\subsection{耦合解释}


TV 距离还有一个非常重要、也非常好用的等价定义：\textbf{两个分布的最优耦合下，两个随机变量不相等的概率}。

设随机变量 $(X,Y)$ 的联合分布以 $P$ 为 $X$ 的边缘、$Q$ 为 $Y$ 的边缘（称为 $P,Q$ 的一个\textbf{耦合}）。则

\[
\boxed{
\Delta_{\mathrm{TV}}(P,Q)
=\min_{\text{coupling }(X,Y)}\Pr[X\ne Y]
}.
\]

\textbf{下界.} 对任意耦合，

\[
\begin{aligned}
\Pr[X\ne Y]
&\ge\Pr[X=x,Y\ne x]\quad\text{对某个 }x\\
&\ge P(x)-Q(x)\quad\text{当 }P(x)>Q(x).
\end{aligned}
\]

对所有 $x$ 取正部并求和，得到 $\Pr[X\ne Y]\ge\Delta_{\mathrm{TV}}(P,Q)$。

\textbf{上界（构造最优耦合）.} 令

\[
m=\sum_x\min\{P(x),Q(x)\}=1-\Delta_{\mathrm{TV}}(P,Q).
\]

构造联合分布：

\begin{itemize}
\item 以 $\min(P(x),Q(x))$ 的概率令 $X=Y=x$（对每个 $x$）；
\item 剩下的概率 $1-m=\Delta_{\mathrm{TV}}(P,Q)$ 放在 $X\ne Y$ 的部分，并按 $P,Q$ 超出最小值的部分分配。
\end{itemize}


这样边缘分别是 $P,Q$，且 $\Pr[X\ne Y]=\Delta_{\mathrm{TV}}(P,Q)$，达到下界。

\textbf{直观.} TV 就是“把两个分布配成一对时，最少有多少概率会不相等”；能配得越像，TV 越小。

\subsection{一个例子：两个 Bernoulli}


设 $P=\mathrm{Bern}(1/3)$，$Q=\mathrm{Bern}(2/3)$。则

\[
\Delta_{\mathrm{TV}}(P,Q)
=\frac12\left(\left|\frac13-\frac23\right|
+\left|\frac23-\frac13\right|\right)
=\frac13.
\]

最优耦合下，两个随机变量不相等的概率可以做到 $1/3$；任何耦合下都至少是 $1/3$。

\subsection{TV 与“不可区分”}


若 $\Delta_{\mathrm{TV}}(P,Q)\le\varepsilon$，则对任何事件 $A$，

\[
|P(A)-Q(A)|\le\varepsilon.
\]

\textbf{证明.} 因为

\[
P(A)-Q(A)=\sum_{x\in A}(P(x)-Q(x)),
\]

取绝对值后不超过 $\sum_x|P(x)-Q(x)|=2\Delta_{\mathrm{TV}}(P,Q)$，所以差不超过 $\varepsilon$。

这说明：如果两个分布 TV 很小，那么任何统计实验都很难区分它们；这就是“统计距离”的名字来源。

\section{KL 散度与条件 KL}


TV 距离是对称的、取值在 $[0,1]$ 的“统计距离”；本节介绍的 \textbf{Kullback–Leibler 散度}（KL divergence）不是对称的，取值可以到 $+\infty$，但它和熵、条件熵、最大似然、假设检验都有很深的关系。

\textbf{本节用到的记号：} $P,Q$ 为同一可数集合上的分布；$P(x),Q(x)$ 为概率质量；$D_{\mathrm{KL}}(P\Vert Q)$ 为 KL 散度；$P_{XY},Q_{XY}$ 为联合分布；$P_{Y\mid X},Q_{Y\mid X}$ 为条件分布。

\subsection{定义}


\[
\boxed{
D_{\mathrm{KL}}(P\Vert Q)
=\sum_xP(x)\log\frac{P(x)}{Q(x)}
}.
\]

约定：

\begin{itemize}
\item $P(x)=0$ 时，该项为 $0$（即 $0\log\frac00=0$）；
\item 若 $Q(x)=0$ 而 $P(x)>0$，则 $D_{\mathrm{KL}}(P\Vert Q)=+\infty$。
\end{itemize}


因此 KL 只在 $P$ 的支撑集包含于 $Q$ 的支撑集时有限；否则为无穷。

\textbf{直观.} $D_{\mathrm{KL}}(P\Vert Q)$ 衡量“用分布 $Q$ 去编码/近似真实分布 $P$”时平均多付出的信息量；$P$ 与 $Q$ 越不像，值越大。

\subsection{基本性质}


\textbf{非负性（Gibbs 不等式）.}

\[
D_{\mathrm{KL}}(P\Vert Q)\ge0,
\]

等号当且仅当 $P=Q$。

\textbf{证明（Jensen）.} 令 $f(t)=t\log t$，它是凸函数。则

\[
\begin{aligned}
D_{\mathrm{KL}}(P\Vert Q)
&=\sum_xQ(x)\frac{P(x)}{Q(x)}
\log\frac{P(x)}{Q(x)}\\
&=\sum_xQ(x)f\left(\frac{P(x)}{Q(x)}\right)\\
&\ge f\left(\sum_xQ(x)\frac{P(x)}{Q(x)}\right)\\
&=f(1)=0.
\end{aligned}
\]

这里用了 $f$ 的凸性和 $\sum_xP(x)=\sum_xQ(x)=1$。等号条件由 $f$ 严格凸给出：$P(x)/Q(x)$ 对所有 $x$ 相同，即 $P=Q$。

\textbf{其他性质.}

\begin{itemize}
\item $D_{\mathrm{KL}}(P\Vert P)=0$。
\item \textbf{不对称}：一般 $D_{\mathrm{KL}}(P\Vert Q)\ne D_{\mathrm{KL}}(Q\Vert P)$。
\item \textbf{不满足三角不等式}，所以不是距离。
\item 对 $Q$ 是凸函数；对 $(P,Q)$ 是联合凸函数。
\end{itemize}


\textbf{与熵的关系.} 若 $H(P)=-\sum_xP(x)\log P(x)$，则

\[
D_{\mathrm{KL}}(P\Vert Q)
=\sum_xP(x)\log\frac1{Q(x)}-H(P).
\]

第一项是“用 $Q$ 编码 $P$ 的交叉熵”，减去 $P$ 自身的熵，就是多出来的部分。

\subsection{二元 KL}


设 $P=\mathrm{Bern}(p)$、$Q=\mathrm{Bern}(q)$，记

\[
d(p\Vert q)
=p\log\frac pq+(1-p)\log\frac{1-p}{1-q}.
\]

（$p,q\in(0,1)$；端点按极限理解。）

\textbf{性质.}

\begin{itemize}
\item $d(p\Vert q)\ge0$，等号当且仅当 $p=q$；
\item $d(p\Vert q)\ne d(q\Vert p)$（一般不对称）；
\item 固定 $q$，$d(p\Vert q)$ 关于 $p$ 是凸函数；
\item 二阶导数
\end{itemize}

\[
  \frac{\partial^2}{\partial p^2}d(p\Vert q)
  =\frac1p+\frac1{1-p}\ge4.
\]
  因此 $d(p\Vert q)$ 在 $p=q$ 处取最小值 $0$，并且有二次下界
\[
  d(p\Vert q)\ge2(p-q)^2.
\]

最后一条可由“$d(q\Vert q)=0$、一阶导在 $q$ 处为 $0$、二阶导至少 $4$”得到：把 $d(\cdot\Vert q)$ 在 $q$ 处 Taylor 展开，二次项系数至少 $2$。

\subsection{条件 KL}


给定 $X$，比较两个条件分布 $P_{Y\mid X}$ 与 $Q_{Y\mid X}$。定义\textbf{条件 KL 散度}为

\[
D_{\mathrm{KL}}(P_{Y\mid X}\Vert Q_{Y\mid X}\mid P_X)
=\sum_xP_X(x)\,
D_{\mathrm{KL}}(P_{Y\mid X=x}\Vert Q_{Y\mid X=x}).
\]

它先固定 $X=x$ 比较两个条件分布，再按 $P_X$ 加权平均。

\subsection{KL 的链式法则}


设

\[
P_{XY}=P_XP_{Y\mid X},\qquad
Q_{XY}=Q_XQ_{Y\mid X}.
\]

则

\[
\boxed{
D_{\mathrm{KL}}(P_{XY}\Vert Q_{XY})
=D_{\mathrm{KL}}(P_X\Vert Q_X)
+D_{\mathrm{KL}}(P_{Y\mid X}\Vert Q_{Y\mid X}\mid P_X)
}.
\]

\textbf{证明.} 把对数拆开：

\[
\begin{aligned}
D_{\mathrm{KL}}(P_{XY}\Vert Q_{XY})
&=\sum_{x,y}P_{XY}(x,y)
\log\frac{P_X(x)P_{Y\mid X}(y\mid x)}
{Q_X(x)Q_{Y\mid X}(y\mid x)}\\
&=\sum_{x,y}P_{XY}(x,y)\log\frac{P_X(x)}{Q_X(x)}\\
&\quad+\sum_{x,y}P_{XY}(x,y)
\log\frac{P_{Y\mid X}(y\mid x)}{Q_{Y\mid X}(y\mid x)}.
\end{aligned}
\]

第一项对 $y$ 求和得

\[
\sum_xP_X(x)\log\frac{P_X(x)}{Q_X(x)}
=D_{\mathrm{KL}}(P_X\Vert Q_X).
\]

第二项按条件 KL 定义就是

\[
\sum_xP_X(x)\sum_yP_{Y\mid X}(y\mid x)
\log\frac{P_{Y\mid X}(y\mid x)}{Q_{Y\mid X}(y\mid x)}
=D_{\mathrm{KL}}(P_{Y\mid X}\Vert Q_{Y\mid X}\mid P_X).
\]

所以链式法则成立。

\textbf{推论（条件 KL 非负）.}

\[
D_{\mathrm{KL}}(P_{Y\mid X}\Vert Q_{Y\mid X}\mid P_X)\ge0,
\]

因为它是若干非负 KL 的加权平均。

\textbf{推论（联合 KL 的分解）.}

\[
D_{\mathrm{KL}}(P_{XY}\Vert Q_{XY})
\ge D_{\mathrm{KL}}(P_X\Vert Q_X),
\]

等号当且仅当条件分布相同（$P_{Y\mid X}=Q_{Y\mid X}$，在 $P_X$ 几乎处处意义下）。

\subsection{与 TV 的对比}


\begin{itemize}
\item \textbf{对称性}：TV 对称；KL 一般不对称。
\item \textbf{三角不等式}：TV 满足；KL 不满足。
\item \textbf{取值范围}：TV 在 $[0,1]$；KL 在 $[0,+\infty]$。
\item \textbf{相同分布}：两者都为 $0$。
\item \textbf{支撑集不匹配}：TV 仍然有限；KL 可能为 $+\infty$。
\item \textbf{常见用途}：TV 用于统计不可区分、耦合；KL 用于编码、假设检验、最大似然。
\end{itemize}


两者也有联系（Pinsker 不等式等），本讲只要求理解定义与链式法则，后续遇到再补。

\subsection{常见错误}


\begin{itemize}
\item 把 KL 当成距离：它不对称，也不满足三角不等式。
\item 忘记约定 $0\log0=0$；以及在 $Q(x)=0,P(x)>0$ 时 KL 为 $+\infty$。
\item 条件 KL 的权重用错：应当用 $P_X(x)$，不是 $Q_X(x)$。
\item 链式法则中第二项漏掉“条件”符号，写成普通的 $D(P_Y\Vert Q_Y)$。
\item 把 $D(P\Vert Q)=0$ 误认为“$P,Q$ 只是很像”；它实际上等价于 $P=Q$。
\end{itemize}


\section*{一页式总览}
\begin{itemize}
\item 熵：$H(X)=\sum_xP_X(x)\log\frac1{P_X(x)}$；$H(X)\ge0$，$H(X)\le\log|\operatorname{supp}(X)|$。
\item 二元熵：$h(p)=p\log\frac1p+(1-p)\log\frac1{1-p}$，在 $p=1/2$ 最大。
\item 联合/条件熵：$H(X,Y)=H(Y)+H(X\mid Y)$；$H(X\mid Y)\le H(X)$，等号 $\iff X,Y$ 独立。
\item 全链式：$H(X_1,\dots,X_n)=\sum_iH(X_i\mid X_1,\dots,X_{i-1})$。
\item TV 距离：$\Delta_{\mathrm{TV}}(P,Q)=\frac12\|P-Q\|_1=\sum_x\max(0,P(x)-Q(x))=1-\sum_x\min(P(x),Q(x))$。
\item 耦合：$\Delta_{\mathrm{TV}}(P,Q)=\min_{\text{coupling}}\Pr[X\ne Y]$。
\item KL：$D_{\mathrm{KL}}(P\Vert Q)=\sum_xP(x)\log\frac{P(x)}{Q(x)}\ge0$，等号 $\iff P=Q$。
\item KL 链式：$D(P_{XY}\Vert Q_{XY})=D(P_X\Vert Q_X)+D(P_{Y\mid X}\Vert Q_{Y\mid X}\mid P_X)$。
\end{itemize}

\section*{参考资料}
\begin{itemize}
\item 课堂手写稿（2026-11-12，两页）与课堂转写。
\item \textbf{[DS]} 邓小铁、王畅。《离散数学与结构》信息论、熵与 KL 散度相关章节。
\item 课程讲义中概率、熵、条件熵、TV 距离与 KL 散度相关部分。
\end{itemize}
\end{document}
